\documentclass[12pt,a4paper]{article}
\usepackage[T2A]{fontenc}
\usepackage[utf8]{inputenc}
\usepackage[ukrainian]{babel}
\usepackage[left=2.5cm,right=2cm,top=2cm,bottom=2.5cm]{geometry}
\usepackage{graphicx}
\usepackage{float}
\usepackage{booktabs}
\usepackage{array}
\usepackage{amsmath}
\usepackage[labelsep=colon]{caption}
\usepackage{enumitem}
\usepackage{fancyvrb}
\usepackage{indentfirst}
\usepackage{tikz}
\usetikzlibrary{positioning,arrows.meta,shapes.geometric}

\graphicspath{{figures/}}
\fvset{fontsize=\small,frame=single,framesep=2mm}
\setlist[enumerate]{leftmargin=*}
\setlength{\parskip}{0.3em}
\setlength{\emergencystretch}{3em}

% Титульна сторінка: #1 — номер лабораторної роботи
\newcommand{\labtitlepage}[1]{%
\begin{titlepage}
\begin{center}
\large
Міністерство освіти і науки України\\
Національний університет ``Львівська політехніка''\\
Кафедра прикладної математики
\vspace*{\fill}

{\LARGE Звіт}\\[0.5em]
{\Large про виконання лабораторної роботи №#1}\\[0.5em]
{\Large з курсу «Основи інженерії даних»}
\vspace*{\fill}
\end{center}
\begin{flushright}
\begin{minipage}{0.4\textwidth}
\large
Виконала:\\
студентка групи ПМІД-11\\
Мазур Д. М.\\[1.5em]
Перевірив:\\
Алєксєєв В. І.
\end{minipage}
\end{flushright}
\vspace*{\fill}
\begin{center}
\large Львів — 2026
\end{center}
\end{titlepage}
\setcounter{page}{2}}

\newcommand{\rsec}[1]{\par\medskip\noindent\textbf{#1}\ }
